% ============================================================================
% Fichier  : partie4/P04_C10_disque_systemes_fichiers.tex : FORN (LPI342)
% Sources : ISO/IEC 27042 ; NIST SP 800-86 ; documentation FAT, NTFS, EXT.
% ============================================================================

\chapter{Forensique Disque et Systèmes de Fichiers}
\label{chap:disque}
\chaptermark{Forensique disque}

\epigraph{%
    « Le système de fichiers est une carte~;\\
    l'enquêteur apprend à lire ses marges autant que son centre. »%
}{Principe du cours}

\niveauChapitre{Intermédiaire à expert}{Partie~III~; administration des
systèmes.}

\begin{boxobjectifs}
\begin{itemize}
    \item Comprendre l'organisation d'un disque et d'une partition.
    \item Distinguer FAT, NTFS et EXT.
    \item Lire les métadonnées de fichiers (dates, propriétés).
    \item Repérer l'espace alloué et non alloué.
    \item Situer les artefacts à examiner.
\end{itemize}
\end{boxobjectifs}

\bigskip

% ============================================================================
\section*{Question directrice : où, dans l'image du disque, chercher la trace de la copie du fichier clients ?}
% ============================================================================

Le disque de M.~NDONGO garde, peut-être, la trace de la copie du fichier
clients sur la clé USB. Ce chapitre apprend à lire l'organisation du disque et
de ses systèmes de fichiers pour y trouver ces traces.

\medskip
\textbf{Plan du chapitre.}
\begin{enumerate}
    \item Disque et partitions (section~\ref{sec:c10-partitions}).
    \item FAT, NTFS, EXT (section~\ref{sec:c10-fs}).
    \item Les métadonnées de fichiers (section~\ref{sec:c10-meta}).
    \item Alloué et non alloué (section~\ref{sec:c10-alloue}).
\end{enumerate}

\separateur

% ============================================================================
\section{Disque et Partitions}
\label{sec:c10-partitions}
% ============================================================================

Un disque se divise en \termecle{partitions}, chacune formatée avec un
\termecle{système de fichiers} qui organise les données. La table des
partitions décrit cette organisation~; une partition cachée ou supprimée est,
en soi, un indice.

% ============================================================================
\section{FAT, NTFS, EXT}
\label{sec:c10-fs}
% ============================================================================

\begin{tableaucours}{L{2.4cm}L{4.4cm}L{5.4cm}}{%
    \tableauHeader{Système} &
    \tableauHeader{Contexte} &
    \tableauHeader{Intérêt forensique}
}
FAT / exFAT & Clés USB, cartes & Simple~; dates de création, modification, accès \\
NTFS & Windows & Riche~: journal, métadonnées, flux, horodatages multiples \\
EXT (2/3/4) & Linux & Inodes, horodatages, journal \\
\end{tableaucours}
\vspace{-0.8em}
\begin{center}{\small\itshape Tableau~10.1. Trois familles de systèmes de fichiers}\end{center}

% ============================================================================
\section{Les Métadonnées de Fichiers}
\label{sec:c10-meta}
% ============================================================================

Chaque fichier porte des \termecle{métadonnées}~: dates de création, de
modification, d'accès, de changement~; nom~; emplacement~; propriétaire. Ces
horodatages, croisés, aident à dater une action (copie, ouverture). On les
interprète avec prudence~: certaines opérations les modifient, d'autres non.

\begin{boxCRO}
\textbf{Lire sans sur-interpréter.} Un horodatage de modification prouve
qu'un fichier a changé, non \emph{qui} l'a changé ni \emph{pourquoi}~:
\cro{0}{+}{0} pour le fait, mais l'interprétation reste prudente. On
distingue le tracé (la date) du déduit (l'intention).
\end{boxCRO}

% ============================================================================
\section{Alloué et Non Alloué}
\label{sec:c10-alloue}
% ============================================================================

\begin{boxscenario}{Opération TSANAGA. La trace de la copie}
\textbf{Constatation.} Sur l'image NTFS du disque, le journal du système de
fichiers et les raccourcis récents montrent qu'un fichier
\texttt{clients.xlsx} a été ouvert puis copié vers un volume amovible le
3~janvier.

\medskip
\textbf{Analyse} (confiance élevée pour la trace, moyenne pour l'auteur). La
copie vers un amovible est établie par les artefacts NTFS. L'identité de la
personne devra être recoupée (session ouverte, autres indices). On date le
fait, on n'affirme pas encore l'auteur.
\end{boxscenario}

\begin{boxmetier}{Analyste forensique}
\textbf{Connaître le système de fichiers.} L'analyste adapte sa lecture au
système~: NTFS offre des horodatages multiples et un journal précieux~; FAT
en offre moins. Savoir ce que chaque système conserve évite les conclusions
erronées.
\end{boxmetier}

\separateur

\begin{boxresume}
\textbf{Ce qu'il faut retenir de ce chapitre~:}
\begin{itemize}
    \item \textbf{Un disque} se divise en partitions et systèmes de fichiers.
    \item \textbf{FAT, NTFS, EXT} conservent des traces de richesse inégale.
    \item \textbf{Les métadonnées} (horodatages) aident à dater une action.
    \item \textbf{On lit avec prudence}~: la date n'est pas l'auteur.
    \item \textbf{L'espace non alloué} garde ce qu'on croyait effacé.
\end{itemize}
\end{boxresume}

\begin{boxinfo}
\textbf{Points de contrôle de l'élève.} Je sais~: décrire disque et
partitions~; comparer FAT, NTFS, EXT~; lire des métadonnées~; distinguer
alloué et non alloué~; dater sans sur-interpréter.
\end{boxinfo}

\bigskip

\begin{boxexo}[title={\ding{45}\quad Exercice~10.1. Partition \niveaubase}]
Qu'est-ce qu'une partition~? Un système de fichiers~?
\end{boxexo}

\begin{boxexo}[title={\ding{45}\quad Exercice~10.2. NTFS \niveaubase}]
Pourquoi NTFS est-il plus riche que FAT pour l'enquête~?
\end{boxexo}

\begin{boxexo}[title={\ding{45}\quad Exercice~10.3. Métadonnées \niveaubase}]
Citez trois horodatages portés par un fichier.
\end{boxexo}

\begin{boxexo}[title={\ding{45}\quad Exercice~10.4. La date et l'auteur \niveaumoyen}]
Pourquoi un horodatage de modification ne désigne-t-il pas l'auteur~?
\end{boxexo}

\begin{boxexo}[title={\ding{45}\quad Exercice~10.5. TSANAGA \niveaumoyen}]
Quels artefacts NTFS montrent la copie de \texttt{clients.xlsx}~?
\end{boxexo}

\begin{boxexo}[title={\ding{45}\quad Exercice~10.6. Partition cachée \niveaumoyen}]
Pourquoi une partition cachée est-elle un indice~?
\end{boxexo}

\begin{boxexo}[title={\ding{45}\quad Exercice~10.7. CRO \niveaumoyen}]
Expliquez pourquoi on lit les horodatages avec prudence.
\end{boxexo}

\begin{boxexo}[title={\ding{45}\quad Exercice~10.8. EXT \niveaumoyen}]
Qu'examine-t-on dans un système de fichiers EXT~?
\end{boxexo}

\begin{boxexo}[title={\ding{45}\quad Exercice~10.9. Plan d'analyse \niveauexpert}]
Rédigez le plan d'analyse disque de l'image NTFS de TSANAGA.
\end{boxexo}

\begin{boxexo}[title={\ding{45}\quad Exercice~10.10. Constat \niveauexpert}]
Rédigez le constat de la copie de \texttt{clients.xlsx} vers l'amovible.
\end{boxexo}

\bigskip

\begin{boxinfo}
\textbf{Pour aller plus loin.}
\begin{itemize}
    \item ISO/IEC 27042~\cite{iso27042}~; NIST SP~800-86~\cite{nist80086}.
\end{itemize}
\end{boxinfo}

\begin{boxinfo}
\textbf{Transition.} Le disque garde aussi ce qu'on a voulu effacer~: le
Chapitre~\ref{chap:artefacts} traite des fichiers supprimés et des artefacts.
\end{boxinfo}
